% CS4246/5446 AY2627 S1 -- Project Proposal
% Body is constrained to ONE page (single-spaced, single-column, >=10pt, 1in margins).
% Title page and references are excluded from that limit per the project guidelines.
\documentclass[10pt,a4paper]{article}

\usepackage[margin=1in]{geometry}
\usepackage{times}
\usepackage[T1]{fontenc}
\usepackage{enumitem}
\usepackage{booktabs}
\usepackage{array}
\usepackage{titlesec}
\usepackage[hidelinks]{hyperref}

% Tight headings so the body fits one page without shrinking the font.
\titleformat{\section}{\normalfont\bfseries}{}{0pt}{}
\titlespacing{\section}{0pt}{3pt plus 1pt minus 1pt}{1pt}
\setlist[itemize]{topsep=1pt,itemsep=0pt,parsep=0pt,leftmargin=*}
\setlength{\parskip}{2pt}
\setlength{\parindent}{0pt}
\pagestyle{empty}

\begin{document}

% ------------------------------------------------------------------ title page
\begin{titlepage}
\centering
\vspace*{2cm}
{\large National University of Singapore \\ School of Computing \par}
\vspace{0.4cm}
{\large CS4246/5446 --- Reinforcement Learning and Sequential Decision Making \\
Semester 1, AY2026/27 \par}
\vspace{1.6cm}
{\LARGE\bfseries Project Proposal \par}
\vspace{0.6cm}
{\Large Learning When \emph{Not} to Trade: Reinforcement Learning for \\[2pt]
Cross-Venue Latency Arbitrage After Transaction Costs \par}
\vspace{1.6cm}
{\large Group X \quad\textbar\quad Application Project \par}
\vspace{1.4cm}

\begin{tabular}{@{}lll@{}}
\toprule
\textbf{Name} & \textbf{Matriculation No.} & \textbf{Email} \\
\midrule
Kerk Tai Heng      & A0XXXXXXX & e1137106@u.nus.edu \\
$\langle$Member 2$\rangle$ & A0XXXXXXX & eXXXXXXX@u.nus.edu \\
$\langle$Member 3$\rangle$ & A0XXXXXXX & eXXXXXXX@u.nus.edu \\
$\langle$Member 4$\rangle$ & A0XXXXXXX & eXXXXXXX@u.nus.edu \\
$\langle$Member 5$\rangle$ & A0XXXXXXX & eXXXXXXX@u.nus.edu \\
$\langle$Member 6$\rangle$ & A0XXXXXXX & eXXXXXXX@u.nus.edu \\
\bottomrule
\end{tabular}

\vspace{0.8cm}
{\small\textit{AI-use disclosure: ChatGPT was used to edit the team's draft for clarity
and concision. The project idea, data, plan, and decisions are the team's work, and the
team remains responsible for the final submission.}\par}

\vfill
{\normalsize Submitted 28 September 2026 \par}
\end{titlepage}

% ------------------------------------------------------------------- body (1p)
\textbf{1. Background and problem.}
Bitcoin perpetual futures trade on both Binance and Hyperliquid, but the two venues do not
incorporate new information at exactly the same time. Binance often moves first, while
Hyperliquid's public order book follows after its next updates; in our recordings, the
slower book has a median information age of about 285\,ms. During this brief window, a move
on Binance can indicate the direction in which Hyperliquid may next move. Latency arbitrage
attempts to trade before that adjustment and exit when the two prices converge.

The visible price gap is not the same as a tradable profit. Its median size is about
1 basis point, while an immediate round trip can cost roughly 7 basis points after taker
fees and the bid--ask spread~\cite{report}. Fees are paid on entry and exit; crossing the
spread means buying at the ask and selling at the bid; execution delay and slippage can
consume more of the remaining edge. A fixed rule that trades whenever the gap exceeds a
threshold therefore loses money after costs. Our question is whether reinforcement
learning (RL) can learn to trade only the unusually strong opportunities and otherwise
remain inactive.

\textbf{2. Data and cost-aware environment.}
We will use roughly 26 million time-aligned, ten-level order-book snapshots collected from
both venues over 30 days. Later days will be reserved for validation and testing so that
the agent never trains on future market conditions. Contiguous periods will be replayed in
a Gymnasium environment~\cite{gym}.

The agent's observation will summarise the cross-venue price gap, each venue's bid--ask
spread, order-book imbalance, short-term volatility, the age of the slower quote, current
inventory, and holding time. Its actions are to stay flat, enter long, enter short, or exit.
The reward is portfolio profit and loss after taker fees, spread and slippage, funding, and
a conservative 150\,ms execution delay. This prevents the agent from treating a visible but
unexecutable price difference as profit.

\textbf{3. RL approach and evaluation.}
We will first implement the fixed-threshold strategy inside the same environment. PPO will
be the main RL method~\cite{ppo}; its conservative policy updates are suitable for noisy
observations and delayed rewards. We will normalise observations and tune reward scaling so
that rare profitable trades are not overwhelmed by the much more common ``do nothing''
steps. If time permits, DQN~\cite{dqn} will provide a value-based comparison for the same
discrete actions.

The baseline and RL agents will be compared on later, unseen days using net return,
drawdown, win rate, number of trades, and trade selectivity. We will repeat the evaluation
with fee assumptions from 0.5 to 3.5 basis points per side and with less favourable spread
and delay assumptions. Success means improving on the fixed rule after identical costs,
not proving guaranteed profitability. A careful estimate of the fee level at which the
edge disappears is also a valid result.

\textbf{4. Scope and deliverables.}
To keep the project feasible, we will study one perpetual-futures market pair and evaluate
only on recorded data. The deliverables are a checked dataset, a replay environment, a
reproducible baseline, an RL policy, and a comparison on held-out periods. Live trading,
multiple assets, and a production execution system are outside the project scope.

\textbf{5. Responsible AI.}
All experiments will be offline and use no live capital. Position and loss limits will be
part of the environment rather than left to the learned policy. We will document data gaps,
cost assumptions, and failed experiments; log the agent's trades and abstentions; and
discuss whether the strategy could disadvantage slower market participants. These measures
support accountability, transparency, and realistic interpretation of the results.

\textbf{6. Division of work.}
\emph{Kerk Tai Heng} --- data preparation and episode generation.
\emph{Member 2} --- replay environment and transaction-cost model.
\emph{Member 3} --- baseline strategy and evaluation tools.
\emph{Member 4} --- PPO training and tuning.
\emph{Member 5} --- held-out evaluation and fee sensitivity tests.
\emph{Member 6} --- Responsible-AI analysis, report integration, and presentation.
All members will review the full pipeline and share testing, writing, and presentation work.

\textbf{7. Milestones and fallback outcomes.}
\textbf{5 Oct:} clean the data, describe the observed price differences, and produce replay
episodes. This can stand alone as a measurement study.
\textbf{19 Oct:} complete the environment, cost model, and baseline. This can stand alone as
a realistic benchmark study.
\textbf{2 Nov:} train the PPO agent and compare it with the baseline on validation data.
This is the minimum complete RL result.
\textbf{9 Nov:} finish held-out evaluation, fee sensitivity tests, Responsible-AI analysis,
and presentation materials. The final report is due 15 Nov.

\textbf{8. Risk assessment.}
\emph{No usable edge after costs} (high): present this as a negative result and estimate the
cost level at which trading becomes unviable.
\emph{RL training is unstable or learns only to abstain} (medium): simplify the state and
reward, compare with a value-based agent, and retain the baseline study as a complete
fallback.
\emph{Results overfit one market period} (medium): use chronological splits, keep the test
period untouched until final evaluation, and report variation across days.
\emph{Simulated fills are too optimistic} (high): use conservative fees, spread, and delay,
then repeat the evaluation under worse assumptions.
\emph{Missing or poor-quality data} (low): exclude affected periods using pre-defined checks
and document the remaining coverage.

% ------------------------------------------------------------------- refs
\newpage
\begin{thebibliography}{9}\small
\bibitem{report} Cross-venue latency arbitrage feasibility study, Binance--Hyperliquid BTC perpetuals, 2026. Internal report supplied to the team.
\bibitem{ppo} Schulman J, Wolski F, Dhariwal P, Radford A, Klimov O. Proximal policy optimization algorithms. arXiv:1707.06347. 2017.
\bibitem{dqn} Mnih V, Kavukcuoglu K, Silver D, et al. Human-level control through deep reinforcement learning. Nature. 2015;518(7540):529--33.
\bibitem{gym} Towers M, Terry JK, Kwiatkowski A, et al. Gymnasium: a standard interface for reinforcement learning environments. arXiv:2407.17032. 2024.
\end{thebibliography}

\end{document}
